\documentclass[11pt]{article}

\usepackage[a4paper,margin=2.5cm]{geometry}
\usepackage{amsmath,amssymb,amsthm}
\usepackage{graphicx}
\usepackage{hyperref}
\usepackage{bm}

\title{
	A Quaternionic Transport Interpretation of\\
	Holt Prime Constellations:\\
	Toward an Arithmetic Double Pendulum
}

\author{
	Thomas Hoffbauer
}

\date{\today}

\newtheorem{definition}{Definition}
\newtheorem{lemma}{Lemma}
\newtheorem{hypothesis}{Hypothesis}

\begin{document}
	
	\maketitle
	
	\begin{abstract}
		We propose a transport-geometric interpretation of the nonconvex prime constellations studied by Fred B. Holt in the cycles of gaps $G(p^\#)$ of Eratosthenes sieve.
		Instead of viewing prime constellations merely as static admissible gap patterns, we interpret them as oriented transport paths on a discrete residue-state geometry derived from the modulo-$12$ prime classes.
		This yields a natural quaternionic holonomy structure, a parent-child coupling operator, and an arithmetic analogue of a coupled double pendulum.
		We formulate first structural lemmas and propose statistical tests for detecting non-isotropic long-range transport structure among admissible prime-gap configurations.
	\end{abstract}
	
	\section{Introduction}
	
	Fred B. Holt recently analyzed admissible nonconvex prime constellations of type
	\[
	(458,3240), \qquad (459,3242),
	\]
	within the cycles of gaps $G(p^\#)$ arising in Eratosthenes sieve \cite{holt}.
	These configurations evolve through primorial cycles via admissible residue transport.
	
	The present note proposes a reinterpretation of these structures as
	\[
	\textit{oriented transport systems}
	\]
	rather than merely admissible gap configurations.
	
	The central idea is that the geometric structure is carried not by individual primes, but by transitions between admissible residue states.
	
	\section{EABC Transport Geometry}
	
	We define four distinguished residue sectors modulo $12$:
	\[
	E \equiv 1,\qquad
	A \equiv 5,\qquad
	B \equiv 7,\qquad
	C \equiv 11
	\pmod{12}.
	\]
	
	A prime-gap constellation
	\[
	s=(g_1,\dots,g_n)
	\]
	induces a transport path
	\[
	x_{k+1}=x_k+g_k \pmod{12}.
	\]
	
	Thus every constellation defines a word
	\[
	W(s)=(x_0,x_1,\dots,x_n)
	\]
	in the discrete state space
	\[
	\mathcal S=\{E,A,B,C\}.
	\]
	
	\section{Transition Matrices}
	
	\begin{definition}
		For a constellation $s$, define its EABC transition matrix
		\[
		M_s(a,b)=\#(a\to b),
		\]
		where $a,b\in\mathcal S$.
	\end{definition}
	
	This converts a prime-gap constellation into a discrete transport operator.
	
	The reversal symmetry of Holt constellations naturally yields:
	\[
	M_{s^{\mathrm{rev}}}=M_s^\top.
	\]
	
	Hence the system possesses a discrete time-reversal symmetry.
	
	\section{Quaternionic Holonomy}
	
	We identify the EABC states with quaternionic units:
	\[
	E\leftrightarrow 1,\qquad
	A\leftrightarrow i,\qquad
	B\leftrightarrow j,\qquad
	C\leftrightarrow k.
	\]
	
	\begin{definition}
		The quaternionic holonomy of a constellation $s$ is
		\[
		Q(s)=\prod_{m=1}^n q(x_m),
		\]
		where $q(x_m)\in\{1,i,j,k\}$.
	\end{definition}
	
	Since quaternion multiplication is noncommutative,
	\[
	ij=k,\qquad
	jk=i,\qquad
	ki=j,
	\]
	the ordering of transitions carries genuine orientation information.
	
	\section{Parent-Child Coupling}
	
	Holt showed that every $(459,3242)$ constellation has the form
	\[
	[2\hat s2],
	\]
	while the $(458,3240)$ constellations arise by removing either the initial or terminal gap $2$ \cite{holt}.
	
	Define:
	\[
	P=[2\hat s2],
	\qquad
	H=[2\hat s],
	\qquad
	T=[\hat s2].
	\]
	
	The removed boundary gap induces the transport
	\[
	C\to E.
	\]
	
	\begin{lemma}[Boundary Coupling Operator]
		For all Holt parent constellations,
		\[
		M(P)=M(H)+R_{CE},
		\]
		where
		\[
		R_{CE}=e_Ce_E^\top
		\]
		is a rank-$1$ boundary transport operator.
	\end{lemma}
	
	This operator acts as the arithmetic coupling joint of the system.
	
	\section{Arithmetic Double Pendulum}
	
	We define the arithmetic double pendulum as:
	\[
	\mathcal D=(M,Q,R),
	\]
	where:
	\begin{itemize}
		\item $M$ is the EABC transport matrix,
		\item $Q$ is the quaternionic holonomy,
		\item $R=R_{CE}$ is the boundary coupling operator.
	\end{itemize}
	
	The corresponding entanglement operators are:
	\[
	\Delta_H=Q(P)^{-1}Q(H),
	\]
	\[
	\Delta_T=Q(P)^{-1}Q(T).
	\]
	
	Preliminary computations indicate:
	\[
	\Delta_H=1,
	\qquad
	\Delta_T=k,
	\]
	suggesting that head modes remain synchronized while tail modes acquire nontrivial chiral holonomy.
	
	\section{Chiral Transport Hypothesis}
	
	\begin{hypothesis}[Chiral Double-Pendulum Hypothesis]
		Among Holt prime constellations,
		\[
		P(\Delta_T\neq 1)
		>
		P(\Delta_H\neq 1).
		\]
		Furthermore, the tail modes preferentially occupy nontrivial quaternionic axes
		\[
		\{\pm i,\pm j,\pm k\}.
		\]
	\end{hypothesis}
	
	This would imply non-isotropic long-range transport structure within the primorial flow.
	
	\section{Discussion}
	
	The proposed framework suggests that:
	\[
	\text{prime constellations}
	\]
	may be interpreted as
	\[
	\text{topologically coupled transport resonators}.
	\]
	
	In this picture:
	\begin{itemize}
		\item admissible residues define local transport channels,
		\item quaternionic holonomies encode orientation,
		\item parent structures induce boundary couplings,
		\item reversal symmetry acts as discrete time reversal,
		\item and long prime-gap constellations behave as coupled arithmetic oscillators.
	\end{itemize}
	
	The essential open question is whether the induced quaternionic transport phases become asymptotically isotropic or exhibit statistically preferred axes.
	
	\section{Future Work}
	
	The next stage is numerical:
	\begin{itemize}
		\item compute holonomy distributions for all Holt constellations,
		\item analyze reversal-pair spectra,
		\item test isotropy of quaternionic axes,
		\item correlate asymptotic populations with transport spectra,
		\item and compare with genuine prime-gap data at large scales.
	\end{itemize}
	
	\begin{thebibliography}{9}
		
		\bibitem{holt}
		Fred B. Holt,
		\textit{On Nonconvex Constellations Among Primes II: $(J,|s|)=(458,3240)$},
		arXiv:2605.19165 (2026).
		
	\end{thebibliography}
	
\end{document}